% SHARD-ID: RC-00-FRONTMATTER
% SHARD-TITLE: Frontmatter, Status Vocabulary, and Scope
% SHARD-SUMMARY: States the aim of the notebook and the two-sided (orbit gas / transfer matrix) picture that organises every section.
% SHARD-SUMMARY: Fixes the status vocabulary and the evidence chain enforced by the lab-book gate.
% SHARD-KEYWORDS: status, scope, claims, provenance, gate
% SHARD-NOTES: none
% SHARD-SCRIPTS: none

\section{Frontmatter, Status Vocabulary, and Scope}

\subsection{Aim}

This lab book records an attempt to formulate the Riemann hypothesis in the
language of tensor networks and operational quantum mechanics: channels, Kraus
operators, transfer matrices, matrix-product states. The aim is a formulation
that captures the mechanisms of the proofs that exist (Weil for curves,
Deligne for varieties over finite fields, Ihara--Bass and the Ramanujan
property for regular graphs) and connects more directly to quantum mechanics
than the Hilbert--P\'olya programme does. Nothing here proves anything about
the Riemann hypothesis. The Markdown notes under \texttt{notes/} are the
exploratory record; this book is their distillation, and every statement here
is tagged.

\subsection{The organising picture}

Every zeta function in sight has two sides. Side A is a nonnegative count of
closed orbits (prime cycles of a graph, points of a curve, prime powers with
weight $\log p$) and its Euler product; it is a product state, the primes do
not talk to each other. Side B is a small non-Hermitian operator whose power
sums reproduce the counts; its eigenvalues are the zeros. The Riemann
hypothesis, in every instance, is the statement that all nontrivial
eigenvalues of side B share one modulus. \Cref{sec:conventions}
fixes the notation and \cref{sec:definitions} the definitions used for both
sides throughout.

\subsection{Status vocabulary}

Every theorem-like environment in this book carries exactly one status tag,
printed at its right margin, equal to the status derived by
\texttt{scripts/labbook\_check.py} from \texttt{db/claims.tsv}:
\begin{itemize}
  \item \texttt{proved}: a written proof exists under \texttt{notes/} and a
  reviewer other than the author has returned \texttt{VERDICT <id>: VALID} in a
  file under \texttt{notes/reviews/}. Both author and reviewers to date are
  Claude models (single family); no second model family has been used.
  \item \texttt{sketched}: written by the author, not yet attacked.
  \item \texttt{numerical}: established by a script whose output is captured
  under \texttt{outputs/}; never more than numerical evidence.
  \item \texttt{cited}: a statement of the literature, backed by a quote in
  \texttt{db/provenance.tsv} that the gate byte-checks against the arXiv TeX
  source under \texttt{refs/src/}.
  \item \texttt{conjectured}, \texttt{open}, \texttt{refuted}, \texttt{assumed}:
  as the words say; a claim that depends on any of the first three prints
  \texttt{unsupported}, and one that depends on an assumption prints with the
  suffix \texttt{-conditional}.
\end{itemize}
Definitions carry no tag; they are \texttt{stipulated} (ours) or
\texttt{quoted} (a source's, with provenance) in \texttt{db/definitions.tsv}.
The derived-status table generated from the claims database is reproduced
here.

% GENERATED by scripts/labbook\_check.py --regen from db/claims.tsv. Do not edit.
{\small\begin{longtable}{llll}\toprule
id & kind & derived status & assumptions \\ \midrule\endhead
\texttt{obs:one-identity} & observation & \texttt{sketched} & \texttt{--} \\
\texttt{obs:rh-finite-permutation} & observation & \texttt{sketched} & \texttt{--} \\
\texttt{obs:rh-small-matrix} & observation & \texttt{sketched} & \texttt{--} \\
\texttt{obs:rh-graph-ramanujan} & observation & \texttt{sketched} & \texttt{--} \\
\texttt{obs:rh-mps-correlation-length} & observation & \texttt{sketched} & \texttt{--} \\
\texttt{obs:rh-flow-pairing} & observation & \texttt{sketched} & \texttt{--} \\
\texttt{obs:hilbert-polya-hermitian} & observation & \texttt{sketched} & \texttt{--} \\
\texttt{prop:scattering-inner} & proposition & \texttt{sketched} & \texttt{--} \\
\texttt{num:scattering-symbol} & numerical & \texttt{numerical} & \texttt{--} \\
\texttt{prop:functional-model-modes} & proposition & \texttt{sketched} & \texttt{--} \\
\texttt{prop:phase-identity} & proposition & \texttt{sketched} & \texttt{--} \\
\texttt{prop:ringnorm-trace} & proposition & \texttt{sketched} & \texttt{--} \\
\texttt{num:ringnorm-3000} & numerical & \texttt{numerical} & \texttt{--} \\
\texttt{prop:bc-mpo} & proposition & \texttt{sketched} & \texttt{--} \\
\texttt{num:bc-mpo} & numerical & \texttt{numerical} & \texttt{--} \\
\texttt{conj:stinespring-jump-form} & conjecture & \texttt{open} & \texttt{--} \\
\texttt{num:weil-lps-ramanujan} & numerical & \texttt{numerical} & \texttt{--} \\
\texttt{num:weil-lps-rh} & numerical & \texttt{numerical} & \texttt{--} \\
\texttt{prop:weil-lps-exactly-ramanujan} & proposition & \texttt{sketched} & \texttt{--} \\
\texttt{conj:weil-lps-hecke} & conjecture & \texttt{open} & \texttt{--} \\
\texttt{prop:as-transfer-matrix} & proposition & \texttt{sketched} & \texttt{--} \\
\texttt{prop:as-unitarity} & proposition & \texttt{sketched} & \texttt{--} \\
\texttt{num:artin-schreier-mps} & numerical & \texttt{numerical} & \texttt{--} \\
\texttt{obs:as-nonquadratic} & observation & \texttt{sketched} & \texttt{--} \\
\texttt{obs:deligne-ingredients} & observation & \texttt{sketched} & \texttt{--} \\
\texttt{obs:abelian-obstruction} & observation & \texttt{sketched} & \texttt{--} \\
\texttt{thm:qihara-general} & theorem & \texttt{proved} & \texttt{--} \\
\texttt{cor:qihara-kraus} & corollary & \texttt{proved} & \texttt{--} \\
\texttt{cor:qihara-unitary} & corollary & \texttt{proved} & \texttt{--} \\
\texttt{cor:qihara-poles} & corollary & \texttt{proved} & \texttt{--} \\
\texttt{obs:rh-iff-ramanujan-channel} & observation & \texttt{sketched} & \texttt{--} \\
\texttt{num:qihara-unitary} & numerical & \texttt{numerical} & \texttt{--} \\
\texttt{num:qihara-general} & numerical & \texttt{numerical} & \texttt{--} \\
\texttt{cit:mo-extended-ihara} & citedfact & \texttt{cited} & \texttt{--} \\
\texttt{cit:mo-adjoint-weights} & citedfact & \texttt{cited} & \texttt{--} \\
\texttt{cit:wf-arbitrary-weights} & citedfact & \texttt{cited} & \texttt{--} \\
\texttt{cit:bc-nonbacktracking} & citedfact & \texttt{cited} & \texttt{--} \\
\texttt{cit:harrow-transfer} & citedfact & \texttt{cited} & \texttt{--} \\
\texttt{cit:iyer-exact-ramanujan} & citedfact & \texttt{cited} & \texttt{--} \\
\texttt{cit:sunada-attribution} & citedfact & \texttt{cited} & \texttt{--} \\
\texttt{obs:prior-art-verdict} & observation & \texttt{sketched} & \texttt{--} \\
\texttt{conj:kraus-ramanujan} & conjecture & \texttt{open} & \texttt{--} \\
\texttt{conj:assemble-over-p} & conjecture & \texttt{open} & \texttt{--} \\
\texttt{obs:prime-by-prime-blind} & observation & \texttt{sketched} & \texttt{--} \\
\texttt{asm:fuchsian-geometry} & assumption & \texttt{assumed} & \texttt{--} \\
\texttt{asm:laplace-spectrum-compact} & assumption & \texttt{assumed} & \texttt{--} \\
\texttt{asm:selberg-zeta-entire} & assumption & \texttt{assumed} & \texttt{--} \\
\texttt{asm:heat-semigroup} & assumption & \texttt{assumed} & \texttt{--} \\
\texttt{asm:zeta-standard} & assumption & \texttt{assumed} & \texttt{--} \\
\texttt{asm:gamma-standard} & assumption & \texttt{assumed} & \texttt{--} \\
\texttt{cit:dz-guillemin} & citedfact & \texttt{cited} & \texttt{--} \\
\texttt{cit:dfg-band} & citedfact & \texttt{cited} & \texttt{--} \\
\texttt{cit:dz-ruelle-selberg} & citedfact & \texttt{cited} & \texttt{--} \\
\texttt{cit:marklof-trace} & citedfact & \texttt{cited} & \texttt{--} \\
\texttt{cit:fgkp-casimir-laplacian} & citedfact & \texttt{cited} & \texttt{--} \\
\texttt{cit:fjs-scattering-det} & citedfact & \texttt{cited} & \texttt{--} \\
\texttt{cit:mv-continuous-term} & citedfact & \texttt{cited} & \texttt{--} \\
\texttt{prop:circle-comb} & proposition & \texttt{proved} & \texttt{--} \\
\texttt{prop:prime-circles-orbital} & proposition & \texttt{proved} & \texttt{--} \\
\texttt{prop:finite-euler-spectrum} & proposition & \texttt{proved-conditional} & \texttt{asm:gamma-standard, asm:zeta-standard} \\
\texttt{prop:lindblad-sum-of-squares} & proposition & \texttt{proved} & \texttt{--} \\
\texttt{prop:casimir-laplacian} & proposition & \texttt{proved} & \texttt{--} \\
\texttt{prop:quantum-lindblad-tensor} & proposition & \texttt{proved} & \texttt{--} \\
\texttt{prop:poincare-jacobian} & proposition & \texttt{proved} & \texttt{--} \\
\texttt{prop:flat-tower} & proposition & \texttt{proved-conditional} & \texttt{asm:fuchsian-geometry} \\
\texttt{prop:tower-divisor-band} & proposition & \texttt{proved-conditional} & \texttt{asm:fuchsian-geometry, asm:laplace-spectrum-compact, asm:selberg-zeta-entire} \\
\texttt{prop:graph-tower-ruelle} & proposition & \texttt{proved-conditional} & \texttt{asm:fuchsian-geometry, asm:selberg-zeta-entire} \\
\texttt{prop:cusp-prime-comb} & proposition & \texttt{proved-conditional} & \texttt{asm:gamma-standard, asm:zeta-standard} \\
\texttt{prop:damping-identity} & proposition & \texttt{proved-conditional} & \texttt{asm:gamma-standard, asm:zeta-standard} \\
\texttt{num:selberg-lindblad} & numerical & \texttt{numerical} & \texttt{--} \\
\texttt{obs:continuous-dictionary} & observation & \texttt{sketched-conditional} & \texttt{asm:fuchsian-geometry, asm:laplace-spectrum-compact, asm:selberg-zeta-entire} \\
\texttt{conj:quantum-lindblad-gap} & conjecture & \texttt{open} & \texttt{--} \\
\bottomrule\end{longtable}}


\subsection{Lockstep rules}

The gate (\texttt{make check}; \texttt{make ci} adds a fresh build) fails
unless: every note under \texttt{notes/} is named by a shard header and every
shard header names an existing note; every evidence script is named by a
shard and has a captured output; every theorem-like environment has one
label, one claims row, and one status tag equal to the derived status; every
definition has one row and appears once; the notation macros are exactly the
generated image of \texttt{db/notation.tsv} with no duplicate symbols; every
provenance quote is found whole and contiguous at its stated locus; and the
build has no undefined references or citations.
